\documentclass[10pt,letterpaper]{article}
\usepackage[utf8]{inputenc}
\usepackage[T1]{fontenc}
\usepackage[spanish]{babel}
\usepackage{amsmath, amssymb}
\usepackage[top=0.8cm, bottom=0.6cm, left=1.2cm, right=1.2cm, includeheadfoot]{geometry}
\usepackage{enumitem}
\usepackage{fancyhdr}
\usepackage{xcolor}
\usepackage{microtype}
\usepackage{array}
\usepackage{tabularx}
\usepackage{helvet}
\usepackage{tcolorbox}
\usepackage{graphicx}
\tcbuselibrary{skins,breakable}
\renewcommand{\familydefault}{\sfdefault}
\setlength{\headheight}{14pt}
\setlength{\parskip}{0pt}
\setlength{\parindent}{0pt}
\linespread{0.9}

\definecolor{applelightgray}{RGB}{180,180,180}

\pagestyle{fancy}
\fancyhf{}
\fancyhead[L]{\textbf{\sffamily Ejercicios de Pr\'actica \textendash\ Probabilidad y Estad\'istica}}
\fancyhead[R]{\small\sffamily Gu\'ia 6 \textendash\ Test de Hip\'otesis}
\fancyfoot[C]{\small\sffamily\thepage}
\renewcommand{\headrulewidth}{0pt}
\fancyhead[C]{\textcolor{applelightgray}{\rule{\textwidth}{0.4pt}}}

\newcommand{\seccion}[1]{%
  \vspace{2pt}%
  \begin{tcolorbox}[enhanced,colback=white,colframe=black,arc=0pt,boxrule=0pt,
    leftrule=2pt,left=6pt,right=6pt,top=1pt,bottom=1pt]
    {\normalsize\bfseries\sffamily #1}
  \end{tcolorbox}%
  \vspace{1pt}%
}

\newcommand{\reglacaja}[1]{%
  \begin{tcolorbox}[enhanced,colback=white,colframe=black,arc=0pt,boxrule=0.5pt,
    left=6pt,right=6pt,top=2pt,bottom=2pt,breakable]
  \small\sffamily #1
  \end{tcolorbox}%
}

\begin{document}
\thispagestyle{empty}

\begin{center}
  {\fontsize{15}{17}\selectfont\bfseries\sffamily Gu\'ia 6 \textendash\ Test de Hip\'otesis\par}
  \vspace{0.05cm}
  {\small\sffamily Hip\'otesis Nula y Alternativa \textperiodcentered\ Estad\'istico de Prueba \textperiodcentered\ Valor Cr\'itico y p-valor \textperiodcentered\ Conclusi\'on en Contexto\par}
  {\rule{3cm}{0.6pt}}
\end{center}

\vspace{0.15cm}

%% ===== PROBLEMA 1 =====
\seccion{Problema 1 \textendash\ Test Unilateral para la Media}

\noindent\sffamily\small Una cl\'inica afirma que el tiempo promedio de espera en el laboratorio cl\'inico es de $30$ minutos. La administraci\'on sospecha que el tiempo real podr\'ia ser \textbf{mayor} por el aumento de demanda. Se seleccion\'o una muestra aleatoria de $64$ pacientes, con tiempo medio de espera de $33{,}2$ minutos y desviaci\'on est\'andar de $8{,}1$ minutos. Use $\alpha=0{,}05$.

\begin{enumerate}[leftmargin=1.2cm, label=\textbf{\alph*)}, itemsep=2pt, topsep=2pt]
  \item Plantee las hip\'otesis nula y alternativa.
  \item Calcule el estad\'istico de prueba.
  \item Determine el valor cr\'itico y la decisi\'on.
  \item Concluya en el contexto del problema.
\end{enumerate}

\reglacaja{%
\textbf{Recuerde:} Para $n\geq30$: $Z=\dfrac{\bar{x}-\mu_0}{s/\sqrt{n}}$. Prueba unilateral derecha: se rechaza $H_0$ si $Z_{calc}>z_{1-\alpha}$ (con $\alpha=0{,}05$: $z=1{,}645$).%
}

%% ===== PROBLEMA 2 =====
\seccion{Problema 2 \textendash\ Test Bilateral para la Media}

\noindent\sffamily\small Una planta envasadora declara que sus envases de suero fisiol\'ogico contienen $500$ ml en promedio. En una inspecci\'on se midi\'o una muestra aleatoria de $49$ envases, obteniendo un contenido medio de $497$ ml con desviaci\'on est\'andar de $7$ ml. Con $\alpha=0{,}05$, ?`existe evidencia de que el contenido promedio difiere de lo declarado?

\begin{enumerate}[leftmargin=1.2cm, label=\textbf{\alph*)}, itemsep=2pt, topsep=2pt]
  \item Plantee las hip\'otesis.
  \item Calcule el estad\'istico de prueba y comp\'arelo con el valor cr\'itico.
  \item Concluya en contexto.
\end{enumerate}

%% ===== PROBLEMA 3 =====
\seccion{Problema 3 \textendash\ Test para una Proporci\'on}

\noindent\sffamily\small Una campa\~na sanitaria se considera exitosa si m\'as del $60\%$ de la poblaci\'on objetivo completa su esquema de vacunaci\'on. En una muestra aleatoria de $200$ personas, $132$ completaron el esquema. Con $\alpha=0{,}05$, ?`hay evidencia de que la campa\~na fue exitosa?

\begin{enumerate}[leftmargin=1.2cm, label=\textbf{\alph*)}, itemsep=2pt, topsep=2pt]
  \item Plantee las hip\'otesis.
  \item Calcule el estad\'istico de prueba.
  \item Decida y concluya en contexto.
\end{enumerate}

\reglacaja{%
\textbf{Recuerde:} Para proporciones: $Z=\dfrac{\hat{p}-p_0}{\sqrt{\dfrac{p_0(1-p_0)}{n}}}$. El error est\'andar se calcula con $p_0$ (valor bajo $H_0$), no con $\hat{p}$.%
}

\newpage

%% ================= RESPUESTAS =================
\seccion{Respuestas}

\small\sffamily
\textbf{Problema 1.}
\textbf{a)} $H_0:\mu=30$ vs $H_1:\mu>30$ (unilateral derecha).
\textbf{b)} $SE=\frac{8{,}1}{\sqrt{64}}=1{,}0125$; $Z_{calc}=\frac{33{,}2-30}{1{,}0125}\approx\mathbf{3{,}16}$.
\textbf{c)} $Z_{crit}=1{,}645$; como $3{,}16>1{,}645$ (p-valor $\approx0{,}0008<0{,}05$), \textbf{se rechaza} $H_0$.
\textbf{d)} Existe evidencia estad\'istica suficiente para afirmar que el tiempo promedio de espera es superior a $30$ minutos.

\medskip
\textbf{Problema 2.}
\textbf{a)} $H_0:\mu=500$ vs $H_1:\mu\neq500$ (bilateral).
\textbf{b)} $SE=\frac{7}{\sqrt{49}}=1$; $Z_{calc}=\frac{497-500}{1}=\mathbf{-3{,}0}$; $|{-3{,}0}|>1{,}96$ (p-valor $\approx0{,}0027$).
\textbf{c)} Se rechaza $H_0$: existe evidencia de que el contenido promedio difiere de los $500$ ml declarados (los envases contienen, en promedio, menos).

\medskip
\textbf{Problema 3.}
\textbf{a)} $H_0:p=0{,}60$ vs $H_1:p>0{,}60$.
\textbf{b)} $\hat{p}=\frac{132}{200}=0{,}66$; $SE_0=\sqrt{\frac{(0{,}6)(0{,}4)}{200}}\approx0{,}0346$; $Z_{calc}=\frac{0{,}66-0{,}60}{0{,}0346}\approx\mathbf{1{,}73}$.
\textbf{c)} $Z_{crit}=1{,}645$; como $1{,}73>1{,}645$ (p-valor $\approx0{,}042<0{,}05$), se rechaza $H_0$: hay evidencia de que m\'as del $60\%$ complet\'o el esquema; la campa\~na fue exitosa.

\end{document}
